%%%%%%%%%%%%%%%%%%%%%%%%%%%%%%%%%%%%%%%%%%%%%%%%%%%%%%%%%%%%%%%%%%%%%%%%%%%
% Función: guía de los formularios de TFM implementados para la EPS-UAH.
%
% Uso: consulta el PDF para preparar los formularios. Normalmente no debes
% modificar este fichero, salvo para mantener la documentación de la plantilla.
%%%%%%%%%%%%%%%%%%%%%%%%%%%%%%%%%%%%%%%%%%%%%%%%%%%%%%%%%%%%%%%%%%%%%%%%%%%

\documentclass[10pt,a4paper]{article}
\usepackage[T1]{fontenc}
\usepackage[utf8]{inputenc}
\usepackage{lmodern}
\usepackage[spanish,es-lcroman]{babel}
\usepackage[margin=2cm]{geometry}
\usepackage{enumitem,tabularx,booktabs}
\usepackage[colorlinks=true,urlcolor=blue,linkcolor=blue]{hyperref}
\usepackage{microtype}
\setlength{\parindent}{0pt}
\setlength{\parskip}{4pt}
\raggedright
\setlength{\emergencystretch}{2em}
\setlist[itemize]{leftmargin=*,nosep,topsep=3pt}
\newcommand{\variable}[1]{\texttt{\textbackslash{}#1}}
\newcommand{\fichero}[1]{\path{#1}}
\title{Guía del papeleo de TFM\\\large Escuela Politécnica Superior --- Universidad de Alcalá}
\author{}
\date{}

\begin{document}
\begin{center}
{\LARGE Guía del papeleo de TFM\par}
\medskip
{\large Escuela Politécnica Superior --- Universidad de Alcalá\par}
\end{center}

\section{Qué encontrarás aquí}

Hemos preparado seis formularios que reutilizan los datos de tu trabajo. Esta guía se centra en la EPS-UAH y distingue las particularidades del Máster Universitario en Ingeniería Electrónica (MUIE). No pretende cubrir todos los másteres de la UAH: comprueba qué normativa y qué formularios corresponden a tu titulación antes de presentar los documentos.

Los convenios, solicitudes y actas se identifican con la normativa de TFM EPS-UAH aprobada en Junta de Centro el 31/10/2025 y por la Comisión de Máster Universitario el 30/01/2026. En los dos formularios compartidos indicamos también las correspondencias con el reglamento UAH modificado el 14/05/2026 y la normativa MUIE-UAH aprobada el 15/06/2026. La correspondencia entre anexos no garantiza que todos los formatos sean intercambiables ni sustituye la comprobación de las versiones oficiales vigentes.

Puedes localizar los formularios mediante este resumen:
\begin{tabularx}{\textwidth}{@{}Xl@{}}
\toprule
Documento & Anexo EPS-UAH \\
\midrule
\hyperref[form:publicacion]{Autorización para publicación en abierto} & XII \\
\hyperref[form:visto-bueno]{Visto bueno de la dirección} & X \\
\hyperref[form:convenio]{Convenio de cooperación educativa} & I \\
\hyperref[form:acta]{Acta de defensa y rúbrica} & XI \\
\hyperref[form:prorroga]{Solicitud de prórroga} & IV \\
\hyperref[form:cambio]{Solicitud de cambio del TFM} & V \\
\bottomrule
\end{tabularx}

\section{Antes de generar los documentos}

Abre \fichero{Config/myconfig.tex} y sustituye los datos de ejemplo por los tuyos. Selecciona tu máster de la EPS-UAH mediante \variable{myDegree}, el idioma mediante \variable{myLanguage} y los títulos mediante \variable{myBookTitleSpanish} y \variable{myBookTitleEnglish}. La selección de titulación proporciona las variables derivadas de universidad, centro y tipo de trabajo; no debes rellenarlas directamente para cada formulario.

El nombre \variable{myAuthorFullName} se obtiene normalmente de \variable{myAuthorName} y \variable{myAuthorSurname}; revisa ambas cuando una ficha te pida el nombre completo. \variable{myBookTitle} utiliza el título del idioma elegido. \variable{myPaperworkDate} toma por defecto la fecha \variable{myThesisDepositDate}, aunque puedes cambiarla. No todos los formularios imprimen esa fecha: las fichas te indican cuáles la usan y cuáles dejan la fecha para su firma o utilizan la fecha del equipo.

\subsection{Género gramatical y singular o plural}

Configura correctamente \variable{myAuthorGender}, \variable{myAcademicTutorGender} y \variable{myCoTutorGender}, con los valores exactos \texttt{male} o \texttt{female}. La plantilla no deduce el género a partir del nombre ni valida otros valores. En el acta también importan \variable{myTribunalPresidentGender}, \variable{myTribunalFirstSpokespersonGender} y \variable{myTribunalSecondSpokespersonGender}. Cada ficha indica las variables que afectan realmente a su texto.

La plantilla adapta, por ejemplo, \emph{autor/autora}, \emph{director/directora}, \emph{el/la}, \emph{D./Dª.}, \emph{Presidente/Presidenta} y \emph{1º/1ª}. Si no tienes cotutor, deja \variable{myCoTutorFullName} vacío: \verb|\newcommand{\myCoTutorFullName}{}|. No escribas \texttt{N/A} como nombre para indicar su ausencia: el nombre vacío controla la omisión de datos y firmas y la selección del singular o plural.

En el visto bueno, un tutor puede aparecer como \emph{tutor} y \emph{da su conformidad}; una tutora, como \emph{tutora} y \emph{da su conformidad}. Si hay cotutor, la plantilla utiliza \emph{tutores} y \emph{dan su conformidad}. Los plurales comunes son actualmente masculinos también si ambas personas son mujeres. Algunos rótulos institucionales permanecen fijos, como \emph{Director/a}; no todos los formularios utilizan todas las variables de género.

\subsection{Casillas configurables y casillas fijas}

No todas las casillas se controlan desde \fichero{Config/myconfig.tex}. Las fichas siguientes indican la variable de cada grupo configurable y señalan las casillas fijas que debes editar en el formulario. Respeta las mayúsculas de los valores: \texttt{Yes}/\texttt{No} para publicación y confidencialidad y \texttt{YES}/\texttt{NO} para la propuesta de Matrícula de Honor. Tras cambiar la configuración, vuelve a compilar y comprueba las marcas en el PDF.

Si una casilla es fija, debes editar el fichero \texttt{.tex} del formulario cuando su selección no corresponda a tu caso: \verb|$\XBox$| imprime una casilla marcada y \verb|$\Box$| una vacía. Intercambia ambos comandos para cambiar la respuesta y deja marcada solo la opción que proceda. Revisa también el texto de la declaración para que no contradiga la marca elegida.

\section{Publicación y autorización para la defensa}

\subsection{Publicación en abierto}\label{form:publicacion}

\fichero{TFM-AutorizacionPublicarAbierto.tex}. Reúne la autorización del estudiante, el director y, si existe, el codirector para la publicación en e-BUAH. Sus referencias normativas son las siguientes:
\begin{itemize}
  \item EPS-UAH: anexo XII; aprobaciones del 31/10/2025 y 30/01/2026.
  \item UAH: anexo 4; reglamento modificado el 14/05/2026.
  \item MUIE-UAH: anexo VIII / anexo 4; normativa aprobada el 15/06/2026.
\end{itemize}
Configura los datos y opciones siguientes:
\begin{itemize}
  \item Estudiante y trabajo: \variable{myAuthorFullName}, \variable{myAuthorDNI}, \variable{myDegree}, \variable{myBookTitleSpanish} y \variable{myBookTitleEnglish} si \variable{myLanguage} es \texttt{english}.
  \item Director: \variable{myAcademicTutorFullName}, \variable{myAcademicTutorDNI} y \variable{myAcademicTutorDepartmentOrInstitution}.
  \item Codirector: \variable{myCoTutorFullName}, \variable{myCoTutorDNI} y \variable{myCoTutorDepartmentOrInstitution}. Su bloque y firma se omiten si su nombre está vacío.
  \item Publicación: \variable{myAuthorizationOpenPublishing}, con \texttt{Yes} o \texttt{No}. Si autorizas, \variable{myAuthorizationOpenPublishingEmbargoMonths} selecciona \texttt{0}, \texttt{6}, \texttt{12}, \texttt{18} o \texttt{24} meses de embargo. No utilices otros valores.
  \item Género: \variable{myAuthorGender}, \variable{myAcademicTutorGender} y \variable{myCoTutorGender}. Adaptan los rótulos y firmas, por ejemplo \emph{la autora}, \emph{la directora} y \emph{el codirector}.
\end{itemize}
Las casillas dependen de \variable{myAuthorizationOpenPublishing}: \texttt{Yes} marca \emph{AUTORIZAN} y deja vacía \emph{NO AUTORIZAN}; además, \variable{myAuthorizationOpenPublishingEmbargoMonths} marca exclusivamente la opción de embargo correspondiente (\texttt{0} significa sin embargo). Con \texttt{No} se marca \emph{NO AUTORIZAN} y quedan vacías todas las opciones de embargo, sea cual sea su valor configurado. Revisa la decisión de publicación y el embargo y recoge las firmas. Este documento no imprime \variable{myPaperworkDate}.

\subsection{Visto bueno de la dirección}\label{form:visto-bueno}

\fichero{TFM-VistoBuenoTutor.tex}. Recoge la conformidad de los directores para depositar y defender el TFM. Corresponde al anexo X de la normativa EPS-UAH (31/10/2025 y 30/01/2026) y al anexo VI de la normativa MUIE-UAH (15/06/2026).

Utiliza \variable{myAcademicTutorFullName}, \variable{myCoTutorFullName}, \variable{myAuthorFullName}, \variable{myBookTitle}, \variable{myPaperworkDate} y \variable{myDegree}. Los tratamientos y el singular o plural dependen de \variable{myAcademicTutorGender}, \variable{myCoTutorGender} y \variable{myAuthorGender}. Las firmas utilizan los nombres de los tutores, no los del tribunal.

Las casillas de conformidad para la defensa son fijas: \emph{Sí} está marcada y \emph{No} está vacía. No existe una variable de configuración para esta decisión; tampoco depende de \variable{myAuthorizationOpenPublishing}. Si la dirección no da su conformidad, debes editar este fichero e intercambiar \verb|$\XBox$| y \verb|$\Box$| en estas dos opciones, no en las de confidencialidad. Ajusta también la frase que afirma que la dirección da su conformidad para que no contradiga la respuesta negativa.

Solo si \variable{myDegree} es \texttt{MUIE}, aparecen las casillas de confidencialidad: \variable{myConfidentialContent} con \texttt{No} marca \emph{NO sujeto a confidencialidad} y deja vacía \emph{SÍ sujeto a confidencialidad}; con \texttt{Yes} ocurre lo contrario. Estas casillas sí son configurables y no debes editar sus comandos directamente. En \texttt{MUIT} y \texttt{MUII} aparece también un espacio para el porcentaje de coincidencias de TURNITIN y la justificación correspondiente, que debes completar manualmente.

\section{Cooperación educativa}

\subsection{Anexo al convenio}\label{form:convenio}

\fichero{TFM-AnexoI-ConvenioCooperacion-EPS-UAH.tex}. Anexo I de la normativa EPS-UAH aprobada el 31/10/2025 y el 30/01/2026. Describe el trabajo realizado en una entidad externa y representa al tutor de la entidad mediante el cotutor.

Revisa estas variables:
\begin{itemize}
  \item Estudiante: \variable{myAuthorFullName}, \variable{myAuthorDNI}, \variable{myAuthorTelephone} y \variable{myAuthorEmail}.
  \item Director académico: \variable{myAcademicTutorFullName}, \variable{myAcademicTutorDNI}, \variable{myAcademicTutorDepartmentOrInstitution} y \variable{myAcademicTutorEmail}.
  \item Entidad y tutor externo: \variable{myCoTutorFullName}, \variable{myCoTutorDepartmentOrInstitution}, \variable{myCoTutorPosition}, \variable{myCoTutorDNI} y \variable{myCoTutorEmail}. Si no hay cotutor, los campos derivados de él muestran \texttt{N/A}.
  \item Trabajo: \variable{myDegree}, \variable{myBookTitleSpanish} y \variable{myThesisAcademicYear}; la titulación proporciona el centro y el tipo de trabajo.
  \item Género: \variable{myAuthorGender} y \variable{myCoTutorGender}, para los tratamientos del estudiante y del tutor externo. Algunos rótulos son fijos.
\end{itemize}
Los restantes campos \texttt{N/A} deben completarse: datos de identificación y dirección de la entidad, teléfonos, calendario, horario y resumen del trabajo. No tienen variables propias en \fichero{myconfig.tex}. La fecha impresa utiliza \variable{today}, la fecha del equipo al compilar, y no \variable{myPaperworkDate}. Recoge las firmas de coordinación, tutor de la entidad y estudiante y comprueba si tu máster exige un modelo distinto.

\section{Evaluación y solicitudes}

\subsection{Acta de defensa y rúbrica}\label{form:acta}

\fichero{TFM-ActaYRubricaDefensa-EPS-UAH.tex}. Anexo XI de la normativa EPS-UAH aprobada el 31/10/2025 y el 30/01/2026. Es el modelo EPS-UAH con su rúbrica: no se convierte automáticamente en el acta específica del MUIE por seleccionar ese máster.

Utiliza \variable{myAuthorFullName}, \variable{myAuthorDNI}, \variable{myDegree}, \variable{myBookTitleSpanish}, \variable{myBookTitleEnglish} si \variable{myLanguage} es \texttt{english}, \variable{myAcademicTutorFullName}, \variable{myCoTutorFullName}, \variable{myDepartment}, \variable{myThesisAcademicYear} y \variable{myTribunalGrade}. Este formulario no suma la nota del tutor: la calificación utilizada es \variable{myTribunalGrade}.

Los firmantes del tribunal son \variable{myTribunalPresident}, \variable{myTribunalFirstSpokesperson} y \variable{myTribunalSecondSpokesperson}. Sus rótulos utilizan \variable{myTribunalPresidentGender}, \variable{myTribunalFirstSpokespersonGender} y \variable{myTribunalSecondSpokespersonGender}, por ejemplo \emph{Presidenta}, \emph{Vocal 1ª} y \emph{Vocal 2º}. Los rótulos del estudiante y de la dirección utilizan \variable{myAuthorGender}, \variable{myAcademicTutorGender} y \variable{myCoTutorGender}. El suplente no aparece en esta acta. Completa las observaciones y la rúbrica y revisa las firmas.

Las casillas de propuesta para Matrícula de Honor dependen de \variable{myTribunalMHProposal}: \texttt{YES} marca \emph{Sí} y deja vacía \emph{No}; \texttt{NO} hace lo contrario. La propuesta no se deduce de \variable{myTribunalGrade}: debes configurar la decisión del tribunal y comprobar que cumple los requisitos aplicables.

\subsection{Solicitud de prórroga}\label{form:prorroga}

\fichero{TFM-AnexoIV-SolicitudProrroga-EPS-UAH.tex}. Anexo IV de la normativa EPS-UAH aprobada el 31/10/2025 y el 30/01/2026. Identifica el trabajo, el estudiante y la dirección que autoriza la solicitud.

Utiliza \variable{myAuthorFullName}, \variable{myBookTitle}, \variable{myThesisAcademicYear}, \variable{myAcademicTutorFullName} y \variable{myCoTutorFullName}. Los tratamientos y rótulos dependen de \variable{myAuthorGender}, \variable{myAcademicTutorGender} y \variable{myCoTutorGender}; el codirector se omite si no está configurado. Recoge las firmas de la dirección y del estudiante. El modelo implementado no incorpora campos específicos de motivos o fecha, ni imprime \variable{myPaperworkDate}; comprueba si necesitas completar esa información según el procedimiento aplicable.

\subsection{Solicitud de cambio del TFM}\label{form:cambio}

\fichero{TFM-AnexoV-SolicitudCambioTFM-EPS-UAH.tex}. Anexo V de la normativa EPS-UAH aprobada el 31/10/2025 y el 30/01/2026. Recoge los datos de la dirección, del estudiante y del trabajo cuyo cambio solicitas.

Utiliza \variable{myAcademicTutorFullName}, \variable{myAcademicTutorDepartmentOrInstitution}, \variable{myCoTutorFullName}, \variable{myCoTutorDepartmentOrInstitution}, \variable{myAuthorFullName}, \variable{myAuthorDNI}, \variable{myAuthorEmail} y \variable{myBookTitle}. Un cotutor vacío produce \texttt{N/A} en sus datos. No usa variables de género para sus rótulos: mantiene \emph{Director/a} y \emph{Codirector/a}. Debes completar la exposición razonada, lo que solicitas y el espacio de firma, marcados como \texttt{N/A}; no hay variables de configuración específicas para esos campos.

\section{Cómo generar los documentos}

No necesitas Make. Abre el formulario que necesites en tu editor de LaTeX y compílalo con PDFLaTeX desde \fichero{PapeleoTFM/}, conservando el resto de la estructura. Si tu editor lo requiere, establece esa carpeta como directorio de trabajo. En Overleaf, selecciona ese fichero como documento principal. Estos formularios no necesitan Biber ni \texttt{makeglossaries}; la rúbrica compartida no se compila por separado.

Si utilizas Make, ejecuta las siguientes órdenes desde \fichero{PapeleoTFM/}; necesitas \texttt{latexmk} además de la distribución de LaTeX:
\begin{minipage}{\linewidth}
\begin{verbatim}
make                         # Todos los formularios
make autorizaciones          # Publicación y visto bueno
make convenios               # Convenio de cooperación
make evaluacion              # Acta y rúbrica
make solicitudes             # Prórroga y cambio de TFM
make TFM-VistoBuenoTutor.pdf
make guia                    # Esta guía, sin generar formularios
make help                    # Opciones disponibles
\end{verbatim}
\end{minipage}
Cada formulario independiente genera un PDF con su mismo nombre. \texttt{make clean} elimina auxiliares y conserva los PDF; \texttt{make cleanall} elimina también los PDF de formularios y de esta guía. También puedes compilar esta guía directamente: no lee tu configuración ni incluye datos personales.

\section{Antes de presentar los documentos}

Aunque la compilación termine sin errores, revisa el resultado:
\begin{itemize}
  \item Comprueba el modelo y la versión que corresponden a tu máster. Esta colección no implementa todos los anexos de las normativas EPS-UAH, UAH o MUIE-UAH.
  \item Verifica nombres, DNI, afiliaciones, títulos, géneros, cotutor, tribunal, curso académico y fechas. Sustituye los datos de ejemplo.
  \item Revisa publicación, embargo y confidencialidad; completa los campos vacíos, \texttt{N/A}, notas, rúbricas y justificaciones. Los nombres y cargos escritos con tratamientos no se corrigen automáticamente.
  \item Obtén las firmas necesarias y utiliza el procedimiento de entrega que te indiquen. La plantilla no firma, presenta ni tramita los documentos por ti.
\end{itemize}
\end{document}
